

\section*{Abstract}
Das T0-Modell bietet zwei mathematisch äquivalente, aber konzeptionell unterschiedliche Berechnungsmethoden für Teilchenmassen: die direkte geometrische Methode und die erweiterte Yukawa-Methode. Beide Ansätze kommen mit der einzigen geometrischen Konstante $\xipar = \frac{4}{3} \times 10^{-4}$ aus; die rationalen Quantenzahlen $(r_i, p_i)$ sind Strukturwahlen, $v$ dient als Umrechnungsanker. Diese Dokumentation enthält zudem Neutrino-Quantenzahlen und die quantenfeldtheoretische Herleitung der $\xi$-Konstante durch EFT-Matching und 1-Loop-Berechnungen. Der Neutrinoteil mit Doppel-$\xi$-Unterdrückung ist nach R109 nicht tragfähig (siehe Nachtrag am Ende); maßgeblich für Neutrinos ist Dok.~340. Für geladene Leptonen und Quarks liegt die durchschnittliche Abweichung unter 1 Prozent; die über zwanzig freien Standardmodell-Parameter der Massen werden dabei auf den einen Parameter $\xi$ und feste rationale Quantenzahlen zurückgeführt.


\section{Einleitung}
\label{sec:introduction}

Die Teilchenphysik steht vor einem fundamentalen Problem: Das Standardmodell mit seinen über zwanzig freien Parametern bietet keine Erklärung für die beobachteten Teilchenmassen. Sie werden als Eingaben gemessen, nicht aus einem tieferen Prinzip abgeleitet. Das T0-Modell setzt hier mit zwei komplementären Berechnungsmethoden an, die mit dem einen Parameter $\xi$ und rationalen Quantenzahlen auskommen; die ursprünglich enthaltene Behandlung der Neutrinomassen ist nach R109 nicht tragfähig (siehe Nachtrag).

\subsection{Das Parameterproblem des Standardmodells}
\label{subsec:parameter_problem}

Trotz seines experimentellen Erfolgs hat das Standardmodell eine theoretische Lücke: Es enthält mehr als 20 freie Parameter, die experimentell bestimmt werden müssen. Dazu gehören:

\begin{itemize}
	\item \textbf{Fermionenmassen}: 9 geladene Lepton- und Quarkmassen
	\item \textbf{Neutrinomassen}: 3 Neutrino-Masseneigenwerte
	\item \textbf{Mischungsparameter}: 4 CKM- und 4 PMNS-Matrixelemente
	\item \textbf{Eichkopplungen}: 3 fundamentale Kopplungskonstanten
	\item \textbf{Higgs-Parameter}: Vakuumerwartungswert und Selbstkopplung
	\item \textbf{QCD-Parameter}: Starke CP-Phase und andere
\end{itemize}

\begin{important}[Parameterreduktion]
	Das T0-Modell führt die über zwanzig freien Parameter des Standardmodells auf \textbf{einen} zurück: die geometrische Konstante $\xipar = \frac{4}{3} \times 10^{-4}$, die aus der Geometrie des dreidimensionalen Raums motiviert ist. Hinzu kommen die rationalen Quantenzahlen $(r_i, p_i)$ als Strukturwahlen und $v$ als Umrechnungsanker. Diese Version enthält zusätzlich Neutrino-Quantenzahlen (nach R109 nicht tragfähig, siehe Nachtrag) sowie die quantenfeldtheoretische Herleitung.
\end{important}

\section{Methodologische Klärung: Etablierung vs. Vorhersage}
\label{sec:methodological_clarification}

\begin{important}[Wissenschaftshistorische Einordnung]
	Das T0-Modell orientiert sich an der Methodik der \textbf{Mustererkennung und systematischen Klassifikation}, wie sie etwa beim Periodensystem (Mendelejew 1869) oder beim Quarkmodell (Gell-Mann 1964) erfolgreich war.
\end{important}

\subsection{Zwei-Phasen-Entwicklung}
\label{subsec:two_phases}

\textbf{Phase 1: Etablierung der Systematik}
\begin{enumerate}
	\item Mustererkennung in bekannten Teilchenmassen (Elektron, Myon, Tau)
	\item Parameterbestimmung aus experimentellen Daten
	\item Etablierung der Quantenzahlenzuordnung
	\item Demonstration der mathematischen Äquivalenz beider Methoden
\end{enumerate}

\textbf{Phase 2: Entfaltung der Vorhersagekraft}
\begin{enumerate}
	\item Extrapolation auf unbekannte Teilchen
	\item Herleitung der Quarksektor aus Leptonmustern
	\item Vorhersage neuer Generationen
	\item Durchführung experimenteller Tests
\end{enumerate}

\subsection{Historische Präzedenzfälle erfolgreicher Musterphysik}
\label{subsec:historical_precedence}

Der Vergleich betrifft die Methode der Mustererkennung, nicht den Stand der Bestätigung:

\begin{table}[H]
	\centering
	\adjustbox{max width=\textwidth}{%
\begin{tabular}{p{3cm}p{4cm}p{4cm}p{3cm}}
		\toprule
		\textbf{Entdeckung} & \textbf{Mustererkennung} & \textbf{Vorhersagen} & \textbf{Bestätigung} \\
		\midrule
		Periodensystem (1869) & Atomgewichte und Eigenschaften & Gallium, Germanium, Scandium & Experimentell bestätigt \\
		Spektrallinien (1885) & Wasserstofflinien & Rydberg-Formel für alle Serien & Quantenmechanik \\
		Quarkmodell (1964) & Hadronenmassen & Achtfacher Weg & QCD-Theorie \\
		\textbf{T0-Modell (2025)} & \textbf{Leptonenmassen} & \textbf{Quarks} & \textbf{Experimentelle Tests} \\
		\bottomrule
	\end{tabular}%
}
	\caption{Historische Präzedenzfälle der Musterphysik}
	\label{tab:historical_precedence}
\end{table}

\section{Von Energiefeldern zu Teilchenmassen}
\label{sec:energy_fields_to_masses}

\subsection{Die fundamentale Herausforderung}
\label{subsec:fundamental_challenge}

Ein zentrales Ergebnis des T0-Modells ist, dass sich Teilchenmassen aus geometrischen Prinzipien berechnen lassen. Während das Standardmodell über 20 freie Parameter zur Beschreibung von Teilchenmassen benötigt, kommt das T0-Modell mit der geometrischen Konstante $\xigeom = \frac{4}{3} \times 10^{-4}$ und rationalen Quantenzahlen aus; die Abweichungen für geladene Leptonen und Quarks liegen unter 1,2 Prozent.

\begin{tcolorbox}[colback=green!5!white,colframe=green!75!black,title=Parameterreduktion bei den Massen]
	\textbf{Parameterreduktion:}
	\begin{itemize}
		\item \textbf{Standardmodell}: 20+ freie Massenparameter (gemessen)
		\item \textbf{T0-Modell}: 1 Parameter $\xi$ plus rationale Quantenzahlen (geometrisch)
		\item \textbf{Experimentelle Genauigkeit}: 99 Prozent durchschnittliche Übereinstimmung (geladene Leptonen und Quarks; Neutrinoteil siehe Nachtrag)
		\item \textbf{Theoretische Grundlage}: Dreidimensionale Raumgeometrie + QFT-Herleitung
	\end{itemize}
\end{tcolorbox}

\subsection{Energiebasiertes Massenkonzept}
\label{subsec:energy_based_mass}

Im T0-Rahmen wird das, was wir traditionell als Masse bezeichnen, als Ausdruck charakteristischer Energieskalen von Feldanregungen gedeutet:

\begin{equation}
	\boxed{m_i \rightarrow E_{\text{char},i} \quad \text{(charakteristische Energie des Teilchentyps } i\text{)}}
	\label{eq:mass_to_energy}
\end{equation}

Diese Umdeutung behandelt Masse und Energie als verschiedene Aspekte derselben Größe.

\section{Zwei komplementäre Berechnungsmethoden}
\label{sec:two_calculation_methods}

Das T0-Modell bietet zwei mathematisch äquivalente, aber konzeptionell unterschiedliche Ansätze zur Berechnung von Teilchenmassen:

\subsection{Methode 1: Direkte geometrische Resonanz}
\label{subsec:direct_geometric_method}

\textbf{Konzeptionelle Grundlage:} Teilchen als Resonanzen im universellen Energiefeld

Die direkte Methode behandelt Teilchen als charakteristische Resonanzmoden des Energiefeldes $\Efield$, analog zu stehenden Wellenmustern:

\begin{equation}
	\text{Teilchen} = \text{Diskrete Resonanzmoden von } \Efield(x,t)
\end{equation}

\textbf{Dreistufiger Berechnungsprozess:}

\textbf{Schritt 1: Geometrische Quantisierung}
\begin{equation}
	\xi_i = \xi_0 \cdot f(n_i, l_i, j_i)
	\label{eq:geometric_quantization}
\end{equation}

Hier ist $f(n,l,j)$ der geometrische Strukturfaktor der Quantenzahlen, keine Frequenz und nicht die Grundwindungszahl $f = 1/\xi$.

wobei:
\begin{align}
	\xi_0 &= \frac{4}{3} \times 10^{-4} \quad \text{(geometrischer Basisp Parameter)} \\
	n_i, l_i, j_i &= \text{Quantenzahlen aus der 3D-Wellengleichung} \\
	f(n_i, l_i, j_i) &= \text{geometrische Funktion aus räumlichen Harmonischen}
\end{align}

\textbf{Schritt 2: Resonanzfrequenzen}
\begin{equation}
	\omega_i = \frac{c^2}{\xi_i \cdot r_{\text{char}}}
	\label{eq:resonance_frequencies}
\end{equation}

In natürlichen Einheiten ($c = 1$):
\begin{equation}
	\omega_i = \frac{1}{\xi_i}
\end{equation}

\textbf{Schritt 3: Massenbestimmung aus Energieerhaltung}
\begin{equation}
	E_{\text{char},i} = \hbar \omega_i = \frac{\hbar}{\xi_i}
	\label{eq:energy_from_frequency}
\end{equation}

In natürlichen Einheiten ($\hbar = 1$):
\begin{equation}
	\boxed{E_{\text{char},i} = \frac{1}{\xi_i}}
	\label{eq:characteristic_energy_direct}
\end{equation}

\subsection{Methode 2: Erweiterte Yukawa-Methode}
\label{subsec:extended_yukawa_method}

\textbf{Konzeptionelle Grundlage:} Brücke zur Standardmodell-Formulierung

Die erweiterte Yukawa-Methode bewahrt die Kompatibilität mit Standardmodell-Berechnungen, während die Yukawa-Kopplungen geometrisch bestimmt statt empirisch angepasst werden:

\begin{equation}
	E_{\text{char},i} = y_i \cdot v
	\label{eq:yukawa_mass_formula}
\end{equation}

wobei $v = 246$ GeV der Higgs-Vakuumerwartungswert ist (SM-Definition: $v\equiv(\sqrt{2}\,G_F)^{-1/2}$; kein unabhängiges Observable, Dok.~351/R112).

\textbf{Geometrische Yukawa-Kopplungen:}
\begin{equation}
	\boxed{y_i = r_i \cdot \left(\frac{4}{3} \times 10^{-4}\right)^{\pi_i}}
	\label{eq:geometric_yukawa}
\end{equation}

\textbf{Generationshierarchie:}
\begin{align}
	\text{1. Generation:} \quad &\pi_i = \frac{3}{2} \quad \text{(Elektron, Up-Quark)} \\
	\text{2. Generation:} \quad &\pi_i = 1 \quad \text{(Myon, Charm-Quark)} \\
	\text{3. Generation:} \quad &\pi_i = \frac{2}{3} \quad \text{(Tau, Top-Quark)}
\end{align}

Die Koeffizienten $r_i$ sind einfache rationale Zahlen, die durch die geometrische Struktur jedes Teilchentyps bestimmt werden.

\section{Quantenfeldtheoretische Herleitung der $\xi$-Konstante}
\label{sec:qft_derivation}

\subsection{EFT-Matching und Yukawa-Kopplung nach EWSB}
\label{subsec:eft_matching}

Nach der elektroschwachen Symmetriebrechung haben wir die Yukawa-Wechselwirkung:

\begin{equation}
	\mathcal{L}_{\text{Yukawa}} \supset -\lambda_h \bar{\psi}\psi H, \quad \text{mit} \quad H = \frac{v + h}{\sqrt{2}}
\end{equation}

Nach EWSB:
\begin{equation}
	\mathcal{L} \supset -m \bar{\psi}\psi - y h \bar{\psi}\psi
\end{equation}

mit den Beziehungen:
\begin{equation}
	m = \frac{\lambda_h v}{\sqrt{2}} \quad \text{und} \quad y = \frac{\lambda_h}{\sqrt{2}}
\end{equation}

Die lokale Massenabhängigkeit vom physikalischen Higgs-Feld $h(x)$ führt zu:

\begin{equation}
	m(h) = m\left(1 + \frac{h}{v}\right) \quad \Rightarrow \quad \partial_\mu m = \frac{m}{v}\partial_\mu h
\end{equation}

\subsection{T0-Operatoren in der effektiven Feldtheorie}
\label{subsec:t0_operators}

In der T0-Theorie treten Operatoren der Form auf:

\begin{equation}
	O_T = \bar{\psi}\gamma^\mu\Gamma_\mu^{(T)}\psi
\end{equation}

mit dem charakteristischen Zeitfeld-Kopplungsterm:
\begin{equation}
	\Gamma_\mu^{(T)} = \frac{\partial_\mu m}{m^2}
\end{equation}

Einsetzen der Higgs-Abhängigkeit:
\begin{equation}
	\Gamma_\mu^{(T)} = \frac{\partial_\mu m}{m^2} = \frac{1}{mv}\partial_\mu h
\end{equation}

Dies legt nahe, dass ein $\partial_\mu h$-gekoppelter Vektorstrom der UV-Ursprung ist.

\subsection{1-Loop-Matching-Berechnung}
\label{subsec:one_loop_matching}

Die 1-Loop-Amplitude für den T0-Vertex ergibt:
\begin{equation}
	F_V(0) = \frac{y^2}{16\pi^2}\left[\frac{1}{2} - \frac{1}{2}\ln\left(\frac{m_h^2}{\mu^2}\right) + \frac{r(r-\ln r-1)}{(r-1)^2}\right]
\end{equation}

Für hierarchische Massen ($m \ll m_h$) dominiert der konstante Term:
\begin{equation}
	F_V(0) \approx \frac{y^2}{32\pi^2}
\end{equation}

\subsection{$\xi$-Formel aus der Higgs-Physik}
\label{subsec:final_xi_formula}

Das EFT-Matching liefert die Beziehung:
\begin{equation}
	\boxed{\xi = \frac{\lambda_h^2 v^2}{16\pi^3 m_h^2}}
\end{equation}

Mit Standard-Higgs-Parametern ($m_h = 125.1$ GeV, $v = 246.22$ GeV, $\lambda_h = m_h^2/(2v^2) \approx 0.129$):
\begin{equation}
	\xi \approx 1.30 \times 10^{-4}
\end{equation}

Dies liegt nahe an der geometrischen Setzung $\xi_0 = \frac{4}{3} \times 10^{-4} \approx 1.333 \times 10^{-4}$ (Abweichung ca. $-2.4$ Prozent). Die Beziehung ist eine Konsistenzprüfung von $\xi$, keine exakte Herleitung. Zur Herkunft des Higgs-Wegs siehe Dok.~385.

\section{Universelle Teilchenmassen-Systematik}
\label{sec:universal_masses}

\subsection{Überarbeitete universelle Tabelle der Fermionen}
\label{subsec:universal_table}

\resizebox{\textwidth}{!}{%
	\begin{tabular}{|l|c|c|c|c|c|l|}
		\hline
		Fermion & Generation & Familie & Spin & $r_f$ & Exponent $p_f$ & SM-Korrespondenz \\
		\hline
		Elektron-Neutrino & 1 & 0 & 1/2 & $4/3$ & $5/2$ & Doppel-$\xi$ \\
		Elektron          & 1 & 0 & 1/2 & $4/3$  & $3/2$ & Leptonenzahl \\
		Myon-Neutrino     & 2 & 1 & 1/2 & $16/5$ & $3$   & Doppel-$\xi$ \\
		Myon              & 2 & 1 & 1/2 & $16/5$ & $1$   & Leptonenzahl \\
		Tau-Neutrino      & 3 & 2 & 1/2 & $25/9$ & $25/9$ & Doppel-$\xi$ \\
		Tau               & 3 & 2 & 1/2 & $25/9$  & $2/3$ & Leptonenzahl \\
		\hline
		Up     & 1 & 0 & 1/2 & $6$          & $3/2$ & Farbe \\
		Down   & 1 & 0 & 1/2 & $\tfrac{25}{2}$ & $3/2$ & Farbe + Isospin \\
		Charm  & 2 & 1 & 1/2 & $2$$^*$          & $2/3$ & Farbe \\
		Strange& 2 & 1 & 1/2 & $\tfrac{26}{9}$ & $1$   & Farbe \\
		Top    & 3 & 2 & 1/2 & $\tfrac{1}{28}$ & $-1/3$ & Farbe \\
		Bottom & 3 & 2 & 1/2 & $\tfrac{3}{2}$  & $1/2$ & Farbe \\
		\hline
	\end{tabular}%
}

\subsubsection{Wicklungszahlen und topologische Zuordnung}

Die rationalen Quantenzahlen $(r_i, p_i)$ sind nicht willkürlich.
Sie folgen aus den Wicklungszahlen $(n_\theta, n_\phi)$ des
4D-Torus $T^4$ gemäß $r_i = n_\phi^2 / n_\theta$ (Dok.~317,
August 2026):

\begin{center}
\begin{tabular}{lcccccc}
\toprule
Lepton & Gen. & $n_\theta$ & $n_\phi$ & $r_i$ & $p_i$ &
  KSAU-Knoten \\
\midrule
$e$          & 1 & 3 & 2 & $4/3$  & $3/2$ & $3_1$ (Trefoil) \\
$\mu$        & 2 & 5 & 4 & $16/5$ & $1$   & $6_3$ (amphicheiral) \\
$\tau$       & 3 & 9 & 5 & $25/9$ & $2/3$ & $7_1$ (Torusknoten) \\
$\nu_e$      & 1 & 3 & 2 & $4/3$  & $5/2$ & --- \\
$\nu_\mu$    & 2 & 5 & 4 & $16/5$ & $3$   & --- \\
$\nu_\tau$   & 3 & 9 & 5 & $25/9$ & $25/9$& --- \\
\bottomrule
\end{tabular}
\end{center}

Die Exponenten-Schrittweite $\Delta p = 1/3$ entspricht der
Modulationsfrequenz der Torus-Windung. Die Knotenzuordnung (KSAU,
Zenodo 2026) ist \textbf{[S]}; die FFGFT-Quantenzahlen sind
\textbf{[K]}. Details in Dok.~317 (KSAU-FFGFT-Synthese).


\footnotetext{* Korrigiert von ursprünglich $8/9$ basierend auf detaillierter numerischer Analyse}

\section{Numerische Rekonstruktion}
\label{sec:complete_numerical_reconstruction}

Die folgende Analyse zeigt die explizite Berechnung aller Fermionen mit beiden Methoden:

\subsection{Grundlagen und experimentelle Eingangsdaten}
\label{subsec:foundations}

\textbf{Fundamentale Konstanten:}
\begin{align}
	\xi_0 = \xi &= \frac{4}{3} \times 10^{-4} = 1.333333333\ldots \times 10^{-4} \\
	v &= 246 \text{ GeV}
\end{align}

\textbf{Experimentelle Massen (PDG-nahe Werte):}
\begin{align}
	m_e^{\text{exp}} &= 0.0005109989461 \text{ GeV} \\
	m_\mu^{\text{exp}} &= 0.1056583745 \text{ GeV} \\
	m_\tau^{\text{exp}} &= 1.77686 \text{ GeV}
\end{align}

\subsection{Geladene Leptonen: Detaillierte Berechnungen}
\label{subsec:charged_leptons_detailed}

\textbf{Elektron-Massenberechnung:}

\textit{Direkte Methode:}
\begin{align}
			\xi_e &= \frac{4}{3} \times 10^{-4} \times f_e(1,0,1/2), \quad f_e = 1.485 \times 10^{7}\ \text{GeV}^{-1} \\
		&= 1980\ \text{GeV}^{-1} \\
		E_{e} &= \frac{1}{\xi_e} = 0.505 \text{ MeV}
\end{align}

\textit{Erweiterte Yukawa-Methode:}
\begin{align}
	r_e &= \frac{4}{3} \\
	y_e &= \frac{4}{3} \times \left(\frac{4}{3} \times 10^{-4}\right)^{3/2} = 2.053 \times 10^{-6} \\
	E_e &= y_e \times 246 \text{ GeV} = 0.505 \text{ MeV}
\end{align}

\textbf{Myon-Massenberechnung:}

\textit{Direkte Methode:}
\begin{align}
			\xi_\mu &= \frac{4}{3} \times 10^{-4} \times f_\mu(2,1,1/2), \quad f_\mu = 7.146 \times 10^{4}\ \text{GeV}^{-1} \\
		&= 9.528\ \text{GeV}^{-1} \\
		E_{\mu} &= \frac{1}{\xi_\mu} = 104.96 \text{ MeV}
\end{align}

\textit{Erweiterte Yukawa-Methode:}
\begin{align}
	y_\mu &= \frac{16}{5} \times \left(\frac{4}{3} \times 10^{-4}\right)^1 = 4.267 \times 10^{-4} \\
	E_\mu &= y_\mu \times 246 \text{ GeV} = 104.96 \text{ MeV}
\end{align}
\textbf{Experiment:} $105.66 \text{ MeV}$ → Abweichung ca. 0.65 Prozent

\textbf{Experiment Elektron:} $0.511 \text{ MeV}$ → Abweichung $-1.18$ Prozent

Die Faktoren $f_i$ der direkten Methode folgen über Gl.~\ref{eq:transformation_bridge} aus den festen Quantenzahlen: $f_e = 1/(r_e\,\xi^{5/2}\,v) = 1.485 \times 10^7$~GeV$^{-1}$ und $f_\mu = 1/(r_\mu\,\xi^{2}\,v) = 7.146 \times 10^4$~GeV$^{-1}$. Sie sind nicht mit den Yukawa-Vorfaktoren $r_i$ gleichzusetzen. Damit liefert die direkte Methode $1/(\xi f_e) = 0.505$~MeV und $1/(\xi f_\mu) = 104.96$~MeV, dieselben Werte wie die Yukawa-Methode.

\subsection{Neutrino-Behandlung}
\label{sec:complete_neutrino_treatment}

\begin{neutrino}[Neutrino-Ansatz]
	Dieser Abschnitt schlägt eine geometrische Behandlung der Neutrinomassen über eine \textbf{Doppel-$\xi$-Unterdrückung} vor. Nach R109 ist dieser Neutrinoteil nicht tragfähig (siehe Nachtrag am Ende); maßgeblich für Neutrinos ist Dok.~340.
\end{neutrino}

\subsection{Neutrino-Quantenzahlen}
\label{subsec:neutrino_quantum_numbers}

Neutrinos folgen derselben Quantenzahlstruktur wie andere Fermionen, jedoch mit einer Modifikation aufgrund ihrer schwachen Wechselwirkungsnatur:

\begin{table}[H]
	\centering
	\adjustbox{max width=\textwidth}{%
\begin{tabular}{lcccc}
		\toprule
		\textbf{Neutrino} & \textbf{n} & \textbf{l} & \textbf{j} & \textbf{Unterdrückung} \\
		\midrule
		$\nu_e$ & 1 & 0 & 1/2 & Doppel-$\xi$ \\
		$\nu_\mu$ & 2 & 1 & 1/2 & Doppel-$\xi$ \\
		$\nu_\tau$ & 3 & 2 & 1/2 & Doppel-$\xi$ \\
		\bottomrule
	\end{tabular}%
}
	\caption{Neutrino-Quantenzahlen mit charakteristischer Doppel-$\xi$-Unterdrückung}
	\label{tab:neutrino_quantum_numbers}
\end{table}

\subsection{Doppel-$\xi$-Unterdrückungsmechanismus}
\label{subsec:double_xi_suppression}

Der Ansatz nimmt an, dass Neutrinos einen zusätzlichen geometrischen Unterdrückungsfaktor erfahren:

\begin{equation}
	f(n_{\nu_i}, l_{\nu_i}, j_{\nu_i}) = f(n_i, l_i, j_i)_{\text{Lepton}} \times \xi
	\label{eq:neutrino_suppression}
\end{equation}

\textbf{Neutrino-Massenberechnungen:}

\textbf{Elektron-Neutrino:}
\begin{align}
	\xi_{\nu_e} &= \frac{4}{3} \times 10^{-4} \times 1 \times \frac{4}{3} \times 10^{-4} = \frac{16}{9} \times 10^{-8} \\
	E_{\nu_e} &= \frac{1}{\xi_{\nu_e}} = 9.1 \text{ meV}
\end{align}

\textbf{Myon-Neutrino:}
\begin{align}
	\xi_{\nu_\mu} &= \frac{4}{3} \times 10^{-4} \times \frac{16}{5} \times \frac{4}{3} \times 10^{-4} = \frac{256}{45} \times 10^{-8} \\
	E_{\nu_\mu} &= \frac{1}{\xi_{\nu_\mu}} = 1.9 \text{ meV}
\end{align}

\textbf{Tau-Neutrino:}
\begin{align}
	\xi_{\nu_\tau} &= \frac{4}{3} \times 10^{-4} \times \frac{25}{9} \times \frac{4}{3} \times 10^{-4} = \frac{400}{81} \times 10^{-8} \\
	E_{\nu_\tau} &= \frac{1}{\xi_{\nu_\tau}} = 18.0 \text{ meV}
\end{align}

\section{Quark-Analyse mit beiden Methoden}
\label{sec:quark_analysis}

\subsection{Explizite Berechnungen der Quarkmassen}
\label{subsec:quark_calculations}

Wir verwenden $\xi=\tfrac{4}{3}\times10^{-4}$ und $v=246\ \mathrm{GeV}$ (SM-Definition aus $G_F$, Dok.~351/R112).
Für die Yukawa-Darstellung:
\[
y_i = r_i\,\xi^{p_i},\qquad m_i^{\text{pred}}=y_i\,v.
\]
Für die direkte geometrische Darstellung:
\[
f_i=\frac{1}{\xi\, m_i^{\text{exp}}},\qquad m_i^{\text{exp}}=\frac{1}{\xi\, f_i}.
\]

\begin{table}[H]
	\centering
		\resizebox{\textwidth}{!}{%
	\begin{tabular}{lcccccc}
		\toprule
		Quark & $p_i$ & $r_i$ (korr.) & $m_i^{\text{pred}}$ & $m_i^{\text{exp}}$ & rel.\ Fehler & Bemerkung\\
		& & & (GeV) & (GeV) & (\%) & \\
		\midrule
		Up     & $3/2$ & $6$        & $2.272\times10^{-3}$ & $2.16\times10^{-3}$ & $+5.21$ & OK \\
		Down   & $3/2$ & $25/2$     & $4.734\times10^{-3}$ & $4.70\times10^{-3}$ & $+0.73$ & OK \\
		Strange& $1$   & $26/9$        & $9.476\times10^{-2}$  & $9.35\times10^{-2}$  & $+1.34$ & OK\\
		Charm  & $2/3$ & $2$      & $1.284\times10^{0}$  & $1.273$             & $+0.87$ & Korrigiert\\
		Bottom & $1/2$ & $3/2$      & $4.261\times10^{0}$   & $4.18$              & $+1.93$ & OK \\
		Top    & $-1/3$& $1/28$     & $1.7198\times10^{2}$  & $172.57$            & $-0.35$ & OK \\
		\bottomrule
	\end{tabular}%
}
	\caption{Yukawa-Vorhersagen mit korrigierten $r_i,p_i$ und Vergleich mit den Massen nach PDG 2024.}
	\label{tab:quark_yukawa_predictions}
\end{table}

\subsection{Korrektur für das Charm-Quark}
\label{subsec:charm_correction}

Der ursprünglich angegebene Wert $r_c=8/9$ reproduziert nicht die Referenzmasse $m_c=1.273\ \mathrm{GeV}$ (PDG 2024). Der erforderliche Wert ist:
\[
r_c^{\text{erforderlich}}=\frac{m_c^{\text{exp}}}{v\,\xi^{2/3}}\approx 1.983 \approx 2.
\]

Daher wurde $r_c \approx 2$ in der korrigierten universellen Tabelle verwendet. Mit $r_c = 2$ exakt ergibt sich $m_c = 1.284$~GeV ($+0.87\,\%$).


\section{Experimenteller Vergleich}
\label{sec:comprehensive_validation}

\subsection{Genauigkeitsanalyse}
\label{subsec:complete_accuracy}

Vergleich der T0-Werte mit den experimentellen Massen:

\begin{table}[H]
	\centering
		\resizebox{\textwidth}{!}{%
	\begin{tabular}{lcccc}
		\toprule
		\textbf{Teilchen} & \textbf{T0-Vorhersage} & \textbf{Experiment} & \textbf{Genauigkeit} & \textbf{Typ} \\
		\midrule
		\multicolumn{5}{c}{\textit{Geladene Leptonen}} \\
		\midrule
		Elektron & 0.505 MeV & 0.511 MeV & 98.82\% & Lepton \\
		Myon & 104.96 MeV & 105.66 MeV & 99.35\% & Lepton \\
		Tau & 1783.4 MeV & 1776.86 MeV & 99.63\% & Lepton \\
		\midrule
		\multicolumn{5}{c}{\textit{Neutrinos}} \\
		\midrule
		$\nu_e$ & 9.1 meV & $< 450$ meV & Kompatibel & Neutrino \\
		$\nu_\mu$ & 1.9 meV & $< 180$ keV & Kompatibel & Neutrino \\
		$\nu_\tau$ & 18.0 meV & $< 18$ MeV & Kompatibel & Neutrino \\
		\midrule
		\multicolumn{5}{c}{\textit{Quarks}} \\
		\midrule
		Up-Quark & 2.272 MeV & 2.16 MeV & 94.79\% & Quark \\
		Down-Quark & 4.734 MeV & 4.70 MeV & 99.27\% & Quark \\
		Strange-Quark & 94.76 MeV & 93.5 MeV & 98.66\% & Quark \\
		Charm-Quark & 1.284 GeV & 1.273 GeV & 99.13\% & Quark \\
		Bottom-Quark & 4.261 GeV & 4.18 GeV & 98.07\% & Quark \\
		Top-Quark & 171.99 GeV & 172.57 GeV & 99.65\% & Quark \\
		\midrule
		\textbf{Durchschnitt} & & & \textbf{98.6\%} & \textbf{Alle Fermionen} \\
		\bottomrule
	\end{tabular}%
}
	\caption{Vergleich der T0-Werte mit experimentellen Massen}
	\label{tab:complete_validation}
\end{table}


\begin{keyresult}[Ergebnis für geladene Leptonen und Quarks]
	Das T0-Modell erreicht 98.6 Prozent durchschnittliche Genauigkeit über die geladenen Leptonen und Quarks, mit dem einen Parameter $\xi$ und festen rationalen Quantenzahlen. Die Neutrino-Zeilen sind nur als Verträglichkeit mit Grenzen angegeben; der Neutrinoteil ist nach R109 nicht tragfähig (siehe Nachtrag).
\end{keyresult}

\section{Vorhersagekraft des etablierten Systems}
\label{sec:predictive_power}

\subsection{Neue Teilchengenerationen}
\label{subsec:new_generations}

Die Extrapolation der etablierten Muster auf eine 4. Generation ergibt:

\textbf{4. Generation (extrapoliert):}
\begin{align}
	n &= 4, \quad \pi_4 = \frac{1}{2}, \quad r_4 \approx 2.0 \\
	m_{\text{4. Gen}} &= r_4 \times \xi^{1/2} \times v \approx 5.7 \text{ GeV}
\end{align}

Die Rechnung ergibt $5.68$~GeV. Ein geladenes Lepton einer 4.~Generation ist unterhalb von etwa $100$~GeV durch LEP ausgeschlossen, und die $Z$-Breite lässt nur drei leichte Neutrinos zu; diese Extrapolation ist daher keine tragfähige Vorhersage.

\subsection{Quarksektor-Extrapolation}
\label{subsec:quark_extrapolation}

Die Leptonenmuster können auf Quarks übertragen werden:

\begin{table}[H]
	\centering

	\adjustbox{max width=\textwidth}{%
\begin{tabular}{lcccc}
		\toprule
		\textbf{Quark} & \textbf{Generation} & \textbf{$r_i$} & \textbf{$\pi_i$} & \textbf{Vorhersage} \\
		\midrule
		Up & 1 & 6 & 3/2 & 2.3 MeV \\
		Down & 1 & 12.5 & 3/2 & 4.7 MeV \\
		Charm & 2 & 2.0 & 2/3 & 1.3 GeV \\
		Strange & 2 & 2.89 & 1 & 95 MeV \\
		Top & 3 & 0.036 & -1/3 & 173 GeV \\
		Bottom & 3 & 1.5 & 1/2 & 4.3 GeV \\
		\bottomrule
	\end{tabular}%
}

	\caption{Quark-Vorhersagen aus etablierten Mustern}
	\label{tab:quark_predictions}
\end{table}

\section{Korrigierte Interpretation der mathematischen Äquivalenz}
\label{sec:corrected_interpretation}

\begin{keypoint}[Bedeutung der Äquivalenz]
	Die mathematische Äquivalenz beider Methoden ist \textbf{per Definition gegeben}, wenn die Parameter ($r_i$ oder $f_i$) aus denselben experimentellen Massen bestimmt werden. Die Äquivalenz ist kein Beweis der Theorie, sondern eine Konsistenzeigenschaft der mathematischen Struktur.
\end{keypoint}

\subsection{Transformationsbeziehung als Brücke}
\label{subsec:transformation_relationship}

Die Beziehung:
\begin{equation}
	f_i = \frac{1}{r_i \, \xi^{\pi_i} \, v \, \xi_0}
	\label{eq:transformation_bridge}
\end{equation}

verbindet beide Methoden mathematisch. Wenn $r_i$ aus experimentellen Massen bestimmt wird, folgt $f_i$ automatisch und umgekehrt.

\begin{table}[H]
	\centering
	\adjustbox{max width=\textwidth}{%
\begin{tabular}{lcccc}
		\toprule
		\textbf{Teilchen} & \textbf{$m^{\text{exp}}$ (GeV)} & \textbf{$r_i$ (Yukawa)} & \textbf{$f_i$ (direkt)} & \textbf{Genauigkeit} \\
		\midrule
		Elektron & 0.000511 & 4/3 & $1.485 \times 10^{7}$ & $98.82\%$ \\
		Myon & 0.10566 & 16/5 & $7.146 \times 10^{4}$ & $99.34\%$ \\
		Tau & 1.77686 & 25/9 & $4.205 \times 10^{3}$ & $99.63\%$ \\
		\midrule
		$\nu_e$ & 9.1 $\times 10^{-6}$ & 1.349 & $8.235 \times 10^{10}$ & Vorhersage \\
		$\nu_\mu$ & 1.9 $\times 10^{-6}$ & 3.221 & $3.947 \times 10^{11}$ & Vorhersage \\
		$\nu_\tau$ & 18.0 $\times 10^{-6}$ & 2.778 & $3.989 \times 10^{10}$ & Vorhersage \\
		\bottomrule
	\end{tabular}%
}
	\caption{Numerische Äquivalenz beider T0-Methoden für alle Leptonen}
	\label{tab:numerical_equivalence_complete}
\end{table}

Bei den geladenen Leptonen stehen die festen Quantenzahlen $4/3$, $16/5$, $25/9$ und die über Gl.~\ref{eq:transformation_bridge} daraus folgenden $f_i$ (in GeV$^{-1}$); beide Methoden ergeben dieselben Massen und damit dieselbe Genauigkeit. Die Neutrinozeilen gehören zum Neutrinoteil, der nicht tragfähig ist (siehe Nachtrag, R109).

\section{Experimentelle Vorhersagen und Präzisionstests}
\label{sec:experimental_predictions}

\subsection{Modifizierte QED-Vertex-Korrekturen}
\label{subsec:qed_corrections}

Die T0-Theorie sagt modifizierte Feynman-Regeln voraus:
\begin{align}
	\text{Zeitfeld-Vertex:} \quad &-i\gamma^\mu\Gamma_\mu^{(T)} = i\gamma^\mu\frac{\partial_\mu m}{m^2} \\
	\text{Modifizierter Fermion-Propagator:} \quad &S_F^{(T0)}(p) = S_F(p) \cdot \left[1 + \frac{\beta}{p^2}\right]
\end{align}

\subsection{Neutrino-Validierung}
\label{subsec:neutrino_validation}

Die T0-Neutrino-Werte sind mit den hier aufgeführten experimentellen Grenzen verträglich (zur Tragfähigkeit des Neutrinoteils siehe Nachtrag, R109):

\begin{table}[H]
	\centering
	\adjustbox{max width=\textwidth}{%
\begin{tabular}{lccc}
		\toprule
		\textbf{Parameter} & \textbf{T0-Vorhersage} & \textbf{Experimentelle Grenze} & \textbf{Status} \\
		\midrule
		$m_{\nu_e}$ & 9.1 meV & $< 450$ meV (KATRIN) & Erfüllt \\
		$m_{\nu_\mu}$ & 1.9 meV & $< 180$ keV (indirekt) & Erfüllt \\
		$m_{\nu_\tau}$ & 18.0 meV & $< 18$ MeV (indirekt) & Erfüllt \\
		$\sum m_\nu$ & 29.8 meV & $< 60$ meV (Kosmologie 2024) & Erfüllt \\
		\bottomrule
	\end{tabular}%
}
	\caption{T0-Neutrino-Vorhersagen vs. experimentelle Einschränkungen}
	\label{tab:neutrino_validation}
\end{table}

\begin{important}[Neutrino-Massen hierarchie]
	Das T0-Modell sagt \textbf{normale Ordnung} voraus: $m_{\nu_\mu} < m_{\nu_e} < m_{\nu_\tau}$, was mit den aktuellen Oszillationsdaten-Präferenzen übereinstimmt.
\end{important}

\section{Methodologische Grundlage}
\label{sec:scientific_legitimacy}

\subsection{Umkehrbarkeit des etablierten Systems}
\label{subsec:reversibility}

Nach der Etablierungsphase kann das T0-System Vorhersagen machen:

\textbf{Etablierte Leptonenmuster:}
\begin{align}
	\text{1. Generation (n=1):} \quad &\pi_i = \frac{3}{2}, \quad r_e = \frac{4}{3} \approx 1.333 \\
	\text{2. Generation (n=2):} \quad &\pi_i = 1, \quad r_\mu = \frac{16}{5} = 3.2 \\
	\text{3. Generation (n=3):} \quad &\pi_i = \frac{2}{3}, \quad r_\tau = \frac{25}{9} \approx 2.778
\end{align}


\subsection{Experimentelle Testbarkeit}
\label{subsec:experimental_testability}

Die T0-Vorhersagen sind experimentell falsifizierbar:

\begin{enumerate}
	\item \textbf{Präzisionsmessungen:} Verfeinerung der $r_i$-Parameter
	\item \textbf{Neutrino-Tests:} Direkte Neutrinomassenmessungen
\end{enumerate}

Methodisch spricht für das T0-Verfahren:

\begin{enumerate}
	\item \textbf{Systematische Struktur:} Alle Parameter folgen erkennbaren Mustern
	\item \textbf{Vorhersagekraft:} Nach der Etablierung werden neue Teilchen vorhersagbar
	\item \textbf{Experimentelle Testbarkeit:} Vorhersagen sind falsifizierbar
	\item \textbf{QFT-Grundlage:} Quantenfeldtheoretische Herleitung der $\xi$-Konstante
	\item \textbf{Historische Vorbilder:} Die Methodik der Musterphysik war mehrfach erfolgreich
\end{enumerate}

\section{Parametersparsamkeit und universelle Struktur}
\label{sec:parameter_free_nature}

\begin{important}[Ein Parameter und feste Quantenzahlen]
	Die Skala der T0-Koeffizienten wird durch $\xi$ bestimmt, das sich auch aus Higgs-Parametern prüfen lässt (Abweichung zu $4/3 \times 10^{-4}$ etwa $-2.4$ Prozent):
	\begin{equation}
		\xi = \frac{\lambda_h^2 v^2}{16\pi^3 m_h^2} \approx 1.30 \times 10^{-4}
	\end{equation}
			Die Beziehung ist eine Konsistenzprüfung, keine Herleitung (siehe Abschnitt~\ref{subsec:final_xi_formula}).
	Zusammen mit den rationalen Quantenzahlen $(r_i, p_i)$ bleibt damit kein kontinuierlich einstellbarer Parameter übrig; $v$ dient als Umrechnungsanker.
\end{important}

\subsection{Universelle Quantenzahlentabelle}
\label{subsec:universal_quantum_table}

\begin{table}[H]
	\centering
	\adjustbox{max width=\textwidth}{%
\begin{tabular}{lcccccc}
		\toprule
		\textbf{Teilchen} & \textbf{n} & \textbf{l} & \textbf{j} & \textbf{$r_i$} & \textbf{$p_i$} & \textbf{Speziell} \\
		\midrule
		\multicolumn{7}{c}{\textit{Geladene Leptonen}} \\
		\midrule
		Elektron & 1 & 0 & 1/2 & 4/3 & 3/2 & -- \\
		Myon & 2 & 1 & 1/2 & 16/5 & 1 & -- \\
		Tau & 3 & 2 & 1/2 & 25/9 & 2/3 & -- \\
		\midrule
		\multicolumn{7}{c}{\textit{Neutrinos}} \\
		\midrule
		$\nu_e$ & 1 & 0 & 1/2 & 4/3 & 5/2 & Doppel-$\xi$ \\
		$\nu_\mu$ & 2 & 1 & 1/2 & 16/5 & 3 & Doppel-$\xi$ \\
		$\nu_\tau$ & 3 & 2 & 1/2 & 25/9 & 25/9 & Doppel-$\xi$ \\
		\midrule
		\multicolumn{7}{c}{\textit{Quarks}} \\
		\midrule
		Up & 1 & 0 & 1/2 & 6 & 3/2 & Farbe \\
		Down & 1 & 0 & 1/2 & 25/2 & 3/2 & Farbe + Isospin \\
		Charm & 2 & 1 & 1/2 & 2 & 2/3 & Farbe \\
		Strange & 2 & 1 & 1/2 & 26/9 & 1 & Farbe \\
		Top & 3 & 2 & 1/2 & 1/28 & -1/3 & Farbe \\
		Bottom & 3 & 2 & 1/2 & 3/2 & 1/2 & Farbe \\
		\bottomrule
	\end{tabular}%
}
	\caption{Universelle Quantenzahlentabelle für alle Fermionen}
	\label{tab:universal_quantum_numbers}
\end{table}

\section{Kritische Bewertung und Grenzen}
\label{sec:critical_evaluation}

\subsection{Theoretische offene Fragen}
\label{subsec:open_questions}

\begin{enumerate}
	\item \textbf{Anzahl der Generationen:} Warum genau drei Generationen?
	\item \textbf{Hierarchieproblem:} Verbindung zwischen verschiedenen Energieskalen
	\item \textbf{CP-Verletzung:} Einbeziehung von CKM- und PMNS-Mischungsmatrizen
\end{enumerate}

\section{Abschließende Bewertung}
\label{sec:concluding_assessment}

\subsection{Wissenschaftlicher Status}
\label{subsec:scientific_status}

Das T0-Modell bietet eine systematische Beschreibung der Teilchenmassen. Die Kombination von:

\begin{itemize}
	\item \textbf{Hoher numerischer Genauigkeit} (98.6 Prozent über geladene Leptonen und Quarks)
	\item \textbf{Parametersparsamkeit} (ein Parameter $\xi$, rationale Quantenzahlen)
	\item \textbf{Breiter Abdeckung} (geladene Leptonen und Quarks; Neutrinos siehe Nachtrag)
	\item \textbf{QFT-Konsistenz} (1-Loop-Herleitung der $\xi$-Konstante)
	\item \textbf{Experimenteller Testbarkeit} (spezifische falsifizierbare Vorhersagen)
\end{itemize}

macht das Modell zu einem prüfbaren Vorschlag.

\subsection{Bedeutung für die fundamentale Physik}
\label{subsec:significance_physics}

Wenn experimentell bestätigt, würde das T0-Modell unser Verständnis der Teilchenphysik deutlich verändern:

\begin{enumerate}
	\item \textbf{Geometrische Interpretation:} Teilchenmassen als Manifestationen der 3D-Raumgeometrie
	\item \textbf{Vereinheitlichung:} Alle Fermionen folgen derselben universellen Struktur
	\item \textbf{Vorhersagekraft:} Neue Teilchen werden aus etablierten Mustern vorhersagbar
	\item \textbf{Einfachheit:} Starke Vereinfachung komplexer Phänomene
\end{enumerate}

Das T0-Modell ist ein Beispiel dafür, dass eine vereinheitlichte Beschreibung nicht unbedingt größere Komplexität erfordert; die Grundstruktur könnte einfach sein.

\begin{thebibliography}{99}
	\bibitem{pascher_t0_energy_2025}
	Pascher, J. (2025). \textit{Das T0-Modell (Planck-referenziert): Eine Reformulierung der Physik}.  
	
	\bibitem{pascher_derivation_2025}
	Pascher, J. (2025). \textit{Feldtheoretische Ableitung des $\beta_T$-Parameters in natürlichen Einheiten ($\hbar = c = 1$)}.  
	
	\bibitem{pascher_qft_2025}
	Pascher, J. (2025). \textit{Vollständige Herleitung der Higgs-Masse und Wilson-Koeffizienten}. T0-Theory Projektdokumentation.
	
	\bibitem{pascher_units_2025}  
	Pascher, J. (2025). \textit{Natürliche Einheitensysteme: Universelle Energiekonversion und fundamentale Längenskala-Hierarchie}.  
	
	\bibitem{katrin_2024}
	KATRIN Collaboration. (2024). \textit{Direct neutrino mass measurement based on 259 days of KATRIN data}. arXiv:2406.13516.
	
	\bibitem{nufit_2024}
	Esteban, I., et al. (2024). \textit{NuFit-6.0: Updated global analysis of three-flavor neutrino oscillations}. J. High Energy Phys. 12, 216.
	
	\bibitem{cosmology_2024}
	Planck Collaboration. (2020). \textit{Planck 2018 results. VI. Cosmological parameters}. Astron. Astrophys. 641, A6.
	
	\bibitem{gell_mann_1964}
	Gell-Mann, M. (1964). \textit{A schematic model of baryons and mesons}. Physics Letters, 8(3), 214--215.
	
	\bibitem{mendeleev_1869}
	Mendelejew, D. (1869). \textit{Über die Beziehungen der Eigenschaften zu den Atomgewichten der Elemente}. Zeitschrift für Chemie, 12, 405--406.
	
	\bibitem{muon_g2_2023}
	Muon g-2 Collaboration. (2023). \textit{Measurement of the positive muon anomalous magnetic moment to 0.20 ppm}. Phys. Rev. Lett. 131, 161802.
	
\end{thebibliography}
% ============================================================
% Nachtrag 7. September 2026 — Marker-Zertifizierung (R105) und
% Einschränkung Neutrinoteil (R109)
% Quelle: Dok. 190-Archiv-2, deklariert 1. September 2026 (R105)
%         und Dok. 190-Archiv-2 (R109)
% Quelldokument bleibt unverändert (append-only, vgl. R50).
% ============================================================

\section*{Nachtrag: Statusklärung (Stand 7.~September 2026)}

\paragraph{R105: Nachträgliche Marker-Zertifizierung (1.~September 2026).}
Dieses Dokument entstand vor Einführung des Marker-Systems (Dok.~A010)
und trägt keine eigenen Marker. Seine Ableitungen sind inhaltlich schlüssig
und als Kettenglieder für das Galois-Programm (Dok.~336, 338, 339) notwendig.

Nachträglicher Status für dieses Dokument:
Wicklungszahlen $r_i$ und Exponenten $p_i$ sowie die Yukawa-Methode
$m_i = r_i\,\xi^{p_i}\,v$ für geladene Leptonen und Quarks: \textbf{[K]}.

Abhängige Dokumente können diese Ergebnisse mit Verweis auf R105 als gesichertes
Kettenglied behandeln, ohne dieses Dokument zu ändern.

\paragraph{R109: Einschränkung --- Neutrinoteil nicht tragfähig \textbf{[X]} (1.~September 2026).}
Die „direkte Methode" $E_i=1/\xi_i$ für Neutrinos ist keine Rechnung im
eigentlichen Sinn: $1/(\xi r_i)$ ergibt für $e:\mu$ das Verhältnis $1:0{,}42$
(inverse Hierarchie, dimensionslos); die daneben stehenden MeV-Werte sind Messwerte.
Die Neutrinomassen $9{,}1/1{,}9/18{,}0$\,meV folgen aus drei unverträglichen
Vorschriften; der Exponent $p_{\nu_\tau}=25/9$ ist $r_\tau$, irrtümlich in die
Exponentenspalte kopiert.

\textbf{Gültig bleibt:} die Yukawa-Methode $m_i=r_i\xi^{p_i}v$ für geladene
Leptonen und Quarks ($<1{,}2\,\%$ Abweichung). Die Zertifizierung \textbf{[K]} aus
R105 gilt für dieses Dokument \emph{nur} für $r_i$, $p_i$ und die Yukawa-Methode ---
nicht für den Neutrinoteil und nicht für die behauptete Äquivalenz beider Methoden.

\textbf{Präzisierung:} Beide Methoden sind über
$f_i=1/(r_i\,\xi^{p_i+1}v)$ (Gl.~\ref{eq:transformation_bridge}) exakt äquivalent und
liefern mit festen $r_i$, $p_i$ dieselben Massen.

\textbf{Maßgeblich für Neutrinos:} Dok.~340 (R103, R108).
